\begin{center}
{\normalsize Universidade de Brasília\par}
{\small Faculdade de Tecnologia — Engenharia de Produção\par}
{\small Engenharia de Produto — PSP6 | Equipe 2 | 2026/2\par}
\vspace{0.6cm}
{\large\bfseries Princípios de solução para a proteção de louças sanitárias\par}
\vspace{0.15cm}
{\normalsize Projeto conceitual do transporte entre distribuidor e obra\par}
\vspace{0.2cm}
{\small Brasília, 7 de novembro de 2026\par}
\end{center}
\vspace{0.25cm}
\section{Delimitação e continuidade do projeto}
A equipe desenvolveu princípios físicos para proteger diretamente uma cuba ou um corpo individual de bacia durante o transporte e a movimentação até a instalação. A primeira matriz de Pugh selecionou a estrutura funcional M1, com pontuação ponderada de +7. M2 foi a referência, com 0, e M3 obteve \( - \)4.

O recorte adotado para M1 é a proteção da louça descoberta. A embalagem original é retirada no distribuidor, onde ocorrem o preparo e a transferência da peça para o dispositivo retornável. O conjunto acompanha a louça até a retirada para instalação. Essa fronteira deve permanecer explícita nas interfaces de operação e na comparação econômica.

A etapa gera quatro opções para cada uma das 28 funções de menor nível. A equipe usa o método TRIZ para tratar dez conflitos do desdobramento da função qualidade (QFD) e organiza três combinações exploratórias. As 23 especificações-meta orientam o dimensionamento e os ensaios das alternativas.

A seleção de M1 é funcional; não representa a escolha de uma concepção física final. C1, C2 e C3 são combinações exploratórias de princípios, ainda sujeitas ao estudo de módulos e interfaces e à comparação posterior. Esta entrega mantém o nível de geração e combinação de princípios do projeto conceitual \citep[cap.~6]{back2008}.

